% SHARD-ID: AQM-29-SPEC-FQ-T-ARITHMETIC-FIELD
% SHARD-TITLE: The Arithmetic Field on \(\operatorname{Spec}\mathbb F_q[t]\)
% SHARD-SUMMARY: Computes the points, residue fields, and Zariski opens of the affine line over \(\mathbb F_q\).
% SHARD-SUMMARY: Builds the closed-point Weyl-Heisenberg net and state presheaf for \(\operatorname{Spec}\mathbb F_q[t]\).
% SHARD-SUMMARY: Shows that regular polynomial profiles are formal infinite-support fields, not quasi-local operators.
% SHARD-KEYWORDS: affine line, finite field, closed point, irreducible polynomial, Weyl-Heisenberg, state presheaf

\section{The Arithmetic Field on \texorpdfstring{\(\operatorname{Spec}\mathbb F_q[t]\)}{Spec Fq[t]}}
\label{sec:spec-fq-t-arithmetic-field}

\subsection{Status and Source Anchors}

This shard instantiates AQM-28 for \(R_1=k[t]\), \(k=\mathbb F_q\).  The
spectrum and \(D(f)\) notation are sourced from the Stacks Project snapshot,
\path{references/algebraic_geometry/stacks_project_algebra.tex}, lines
2972--2983 and 3104--3117.  Residue fields are sourced from lines 250--274
and 3310--3350.  The polynomial-ring classification below is proved directly
from the Euclidean property of \(k[t]\).  The Weyl layer is the finite
closed-point construction of AQM-28.

\subsection{Lamport Dependency Skeleton}

\[
\begin{array}{rcl}
D1&:&R_1=k[t],\quad X_1=\Spec R_1.\\
L1&:&X_1=\{\eta=(0)\}\cup\{x_\pi=(\pi):\pi
      \text{ monic irreducible}\}.\\
L2&:&\kappa(\eta)=k(t),\quad
      \kappa(x_\pi)=k[t]/(\pi),\quad
      |\kappa(x_\pi)|=q^{\deg\pi}.\\
L3&:&D(g)=\emptyset\text{ if }g=0,\text{ and for }g\ne0
      \text{ it omits exactly the closed points dividing }g.\\
D2&:&E_1(U)=\bigoplus_{x_\pi\in U}\kappa(x_\pi)^2,\quad
      \mathcal A_1(U)=\varinjlim_{S\Subset C(U)}\mathcal A(S).\\
T1&:&\mathcal A_1(U)\text{ is the open-local Weyl algebra on closed
      points of }U.\\
T2&:&\operatorname{State}(\mathcal A_1(U))\text{ is the inverse limit of
      finite-volume density-matrix states.}\\
T3&:&\text{Nonzero polynomial field profiles have infinite closed-point
      support.}
\end{array}
\]

\subsection{Points of the Affine Line L1--L2}

Let \(\mathcal P_k\) be the set of monic irreducible polynomials in \(k[t]\).
The zero ideal is prime because \(k[t]\) is a domain; write
\[
  \eta=(0)
\]
for the generic point.

\paragraph{Lemma \(L1\).}
The remaining prime ideals are exactly
\[
  x_\pi=(\pi),\qquad \pi\in\mathcal P_k.
\]

\emph{Proof.}
Since \(k[t]\) is Euclidean, every nonzero ideal has a unique monic generator.
Let \(\mathfrak p\ne(0)\) be prime and write \(\mathfrak p=(g)\) with \(g\)
monic.  If \(g=ab\) with nonconstant \(a,b\), then \(ab\in\mathfrak p\), so
primality gives \(a\in\mathfrak p\) or \(b\in\mathfrak p\), contradicting
minimal degree of the monic generator \(g\).  Thus \(g\) is irreducible.
Conversely, if \(\pi\) is irreducible, then \(k[t]/(\pi)\) is a field, so
\((\pi)\) is maximal and therefore prime.

\paragraph{Lemma \(L2\).}
\[
  \kappa(\eta)=k(t),\qquad
  \kappa(x_\pi)=k[t]/(\pi).
\]
If \(d=\deg\pi\), then \(\kappa(x_\pi)\) has \(q^d\) elements.  Thus the
closed point \(x_\pi\) carries one \(q^d\)-level finite-field qudit.

\emph{Proof.}
The residue field at the zero prime is the fraction field of \(k[t]\), namely
\(k(t)\).  For the maximal ideal \((\pi)\), localization does not change the
quotient residue field, so \(\kappa(x_\pi)=k[t]/(\pi)\).  The residue classes
of \(1,t,\ldots,t^{d-1}\) form a \(k\)-basis, hence the quotient has \(q^d\)
elements.

The generic point \(\eta\) has infinite residue field \(k(t)\).  It is not a
finite Weyl site under the closed-point convention of AQM-28.

\subsection{Zariski Opens L3}

Let \(g\in k[t]\).  If \(g=0\), then \(D(g)=\emptyset\).  If
\[
  g=u\prod_{i=1}^r\pi_i^{n_i}
\]
with \(u\in k^\times\) and distinct monic irreducibles \(\pi_i\), then
\[
  D(g)=\{\eta\}\cup
  \{x_\pi:\pi\in\mathcal P_k,\ \pi\notin\{\pi_1,\ldots,\pi_r\}\}.
\]

\emph{Derivation.}
The generic prime \((0)\) contains no nonzero \(g\), so \(\eta\in D(g)\).
The closed prime \((\pi)\) contains \(g\) iff \(\pi\) divides \(g\).  Therefore
\(D(g)\) is exactly the displayed set.  In particular every nonempty open
contains \(\eta\), and its closed-point support is cofinite in
\(\mathcal P_k\).

For a closed point \(x_\pi\),
\[
  \{x_\pi\}=V(\pi),\qquad
  X_1\setminus\{x_\pi\}=D(\pi).
\]
Thus closed singletons are not open, but their complements are standard opens.

\subsection{Weyl-Heisenberg Field Net D2--T1}

For an open \(U\subset X_1\), define
\[
  C_1(U)=U\cap\{x_\pi:\pi\in\mathcal P_k\}.
\]
The affine-line field-label group is
\[
  E_1(U)=\bigoplus_{x_\pi\in C_1(U)}\kappa(x_\pi)^2.
\]
For \(D(g)\) with \(g\ne0\),
\[
  E_1(D(g))=
  \bigoplus_{\pi\nmid g}\bigl(k[t]/(\pi)\bigr)^2.
\]
The corresponding algebra is
\[
  \mathcal A_1(D(g))=
  \varinjlim_{S\Subset\{\pi:\pi\nmid g\}}
  \bigotimes_{\pi\in S}M_{q^{\deg\pi}}(\mathbb C).
\]
For \(g=1\), this is the global algebra
\[
  \mathcal A_1(X_1)=
  \varinjlim_{S\Subset\mathcal P_k}
  \bigotimes_{\pi\in S}M_{q^{\deg\pi}}(\mathbb C).
\]

\paragraph{Theorem \(T1\).}
The assignment \(U\mapsto\mathcal A_1(U)\) is the open-local
Weyl-Heisenberg net on \(\Spec k[t]\).

\emph{Proof.}
This is AQM-28 applied to \(R_1=k[t]\).  If \(U\subset V\), then
\(C_1(U)\subset C_1(V)\), so finite closed supports in \(U\) are finite closed
supports in \(V\).  The algebra inclusion tensors by identities on the extra
closed points.  Additivity for open covers follows because every Weyl label
has finite closed support.

At a single closed point \(x_\pi\), let \(K_\pi=k[t]/(\pi)\).  For
\(a,b\in K_\pi\), write
\[
  X_\pi(a)=W_\pi(a,0),\qquad Z_\pi(b)=W_\pi(0,b).
\]
Then
\[
  Z_\pi(b)X_\pi(a)
  =
  \psi_\pi(ab)X_\pi(a)Z_\pi(b),
\]
while operators at distinct closed points commute.

\subsection{States T2}

For every open \(U\),
\[
  \operatorname{State}(\mathcal A_1(U))
  \cong
  \varprojlim_{S\Subset C_1(U)}
  \operatorname{State}\!\left(
  \bigotimes_{x_\pi\in S}M_{q^{\deg\pi}}(\mathbb C)\right).
\]
The restriction maps are partial traces over omitted closed points.  The
product tracial state is obtained by taking the normalized trace at every
closed point; it sends every nontrivial finite-support Weyl operator to \(0\).

The state presheaf is not separated for the Zariski topology.  Let
\[
  x_0=(t),\qquad x_1=(t-1),\qquad
  U=D(t),\quad V=D(t-1).
\]
Then \(U\cup V=X_1\).  On the two \(q\)-dimensional factors \(x_0,x_1\), take
the maximally entangled vector
\[
  |\Phi\rangle=q^{-1/2}\sum_{a\in k}|a,a\rangle.
\]
Let \(\rho=|\Phi\rangle\langle\Phi|\), and let
\(\sigma=(I_q/q)\otimes(I_q/q)\).  Tensor both with the product tracial state
on all other closed points.  The two global states restrict equally to \(U\)
and to \(V\), because each restriction traces out one of \(x_0,x_1\).  They
differ globally, since \(\rho\ne\sigma\) on the two-site subalgebra.

\subsection{Polynomial Field Profiles T3}

A pair of polynomials \(a,b\in k[t]\) defines a closed-point profile
\[
  e_{a,b}(x_\pi)=(a\bmod\pi,\ b\bmod\pi).
\]
There are infinitely many monic irreducibles in \(k[t]\): if
\(\pi_1,\ldots,\pi_n\) were all of them, then
\(P=\pi_1\cdots\pi_n\) is nonconstant and \(P+1\) has an irreducible divisor
which divides none of the \(\pi_i\), a contradiction.
If \(a\ne0\), then \(a\bmod\pi=0\) only for the finitely many irreducible
divisors of \(a\).  Hence \(a\) is nonzero at infinitely many closed points.
The same holds for \(b\).  Therefore
\[
  e_{a,b}\in E_1(X_1)
  \quad\Longleftrightarrow\quad
  a=b=0.
\]

Thus the coordinate function \(t\), and every nonzero polynomial profile, is a
formal infinite-support field profile.  Its finite restrictions
\[
  e_{a,b}|_S\in\bigoplus_{x_\pi\in S}\kappa(x_\pi)^2
\]
give finite-volume Weyl operators, but the full profile does not define an
element of the algebraic quasi-local algebra.

\subsection{Conclusion}

\(\Spec k[t]\) produces a countable closed-point array of finite-field qudits:
one \(q^d\)-level qudit for each monic irreducible polynomial of degree \(d\).
Nonempty Zariski opens remove finitely many closed points and always contain
the generic point, but the generic point has residue field \(k(t)\) and is not
a finite Weyl site in the current convention.  The quasi-local observables are
finite closed-support Weyl operators; polynomial fields are formal
infinite-volume profiles until a completion convention is added.
